
\documentclass[12pt]{article}
\usepackage[utf8]{inputenc}
\usepackage[brazil]{babel}
\usepackage{graphicx}
\usepackage{float}
\usepackage{epsfig}
\usepackage{mathptmx} %para Times Roman
\usepackage{setspace} \setstretch{1.15} %espaçamento 1.15
\usepackage{geometry} %para definir a geometria da página
\usepackage{booktabs} %pacote para tabelas
\usepackage{caption} %Legendas das figuras
\geometry{left=3cm,right=2cm, top=5cm, bottom=2cm} %margens da página
\usepackage{fancyhdr}
\setlength{\parindent}{1.5cm}
\usepackage{indentfirst}
\usepackage{background}
\backgroundsetup{contents={\includegraphics[scale=1.3]{logo.png}},
	scale=0.8,placement=top,opacity=0.5,position={10cm,1.1cm}, vshift=15}
\pagestyle{empty}%retirando numeração
\cfoot{\textbf{\footnotesize CTJ 219 Microprocessadores e Microcontroladores.} Janaúba, MG, Turma A, Período 2021/1.}
\rfoot{\thepage}
\pagestyle{fancy}
\everymath{\displaystyle}
\newtheorem{exem}{Exemplo}


\begin{document}
\pagenumbering{arabic}{
\begin{center}
    \textbf{TÍTULO EM TIMES 12, CENTRALIZADO, NEGRITO E MAIÚSCULAS: SUBTÍTULO TAMBÉM TODO EM TIMES 12, NEGRITO E MAÍUSCULAS, SENDO A EXTENSÃO MÁXIMA DO CONJUNTO DE TRÊS LINHAS }
\end{center}	

\begin{center}
{	\textbf{SOBRENOME, Nome do Autor};
    SOBRENOME, Nome do Autor;
    SOBRENOME, Nome do Autor\\

}
\end{center}

\begin{center}
    Formação ou titulação do autor, instituição de ensino (abreviatura do nome da instituição), cidade, estado, e-mail, o nome do apresentador deverá estar em negrito.\\
    Formação ou titulação do autor, instituição de ensino (abreviatura do nome da instituição), cidade, estado, e-mail.\\
\end{center}


{\setlength{\parindent}{0cm}

\renewcommand{\abstractname}{{\bf RESUMO}}
\begin{abstract}
    {\small
    \noindent Deve ser elaborado em fonte Times 11 justificado, em parágrafo único, com recuos de 1 cm à direita e à esquerda, mas sem recuo na primeira linha, com espaço entrelinhas simples, com o título em negrito. Deve usar o estilo “Resumo” presente no documento. Deve conter no máximo de 150 palavras.
}
\end{abstract}

 
\vspace{2ex}

\textbf{PALAVRAS-CHAVE}: No mínimo 3 (três) e no máximo 5 (cinco) palavras-chave. 

}
\section*{\large INTRODUÇÃO}

    O artigo deve ser realizado para uma página tamanho padrão A4. no formato retrato. A página do artigo deverá ter margem superior igual a 3,0 cm e as demais igual a 2,0 cm. Serão aceitos apenas artigos inéditos. \\
    O texto do artigo conterá as seguintes seções: INTRODUÇÃO, METODOLOGIA, FUNDAMENTAÇÃO TEÓRICA E OU DISCUSSÕES, CONSIDERAÇÕES FINAIS e REFERÊNCIAS. O espaçamento entre linhas será simples, fonte Times 12 e o parágrafo será justificado.\textbf{O RESUMO EXPANDIDO DEVERÁ TER DE DUAS A QUATRO PÁGINAS}.\\
    A introdução deve conter uma referência ao assunto a ser desenvolvido no resumo expandido, bem como as linhas gerais que serão desenvolvidas no corpo do texto. Tal seção não admitirá subdivisões. A Introdução deverá conter o(s) objetivo(s) do estudo apresentado. \\

\section*{\large METODOLOGIA}

    Na METODOLOGIA será explicitado o tipo de estudo, local, população (caso for pesquisa de campo), período, técnica e análise dos dados, bem como as normas éticas seguidas que foram utilizadas no caso de a pesquisa ser com seres humanos, enfim todos os métodos utilizados para a realização do trabalho. \\


\section*{\large RESULTADOS E DISCUSSÕES}

    O TÍTULO DA SEÇÃO DE DESENVOLVIMENTO deverá ser escrito em negrito e maiúsculas, após o título um espaço (uma linha em branco) separará o título da seção do texto correspondente a ela.\\ 
    Não haverá espaços (linhas em branco) entre os parágrafos. A primeira linha de cada um dos parágrafos terá um recuo de 1,5cm. Nesse item não deverá conter citações diretas. \\


\section*{\large CONSIDERAÇÕES FINAIS}

    Favor seguir as normas de diagramação aqui expostas, usando este exemplo como base para o seu texto. \\
    Nas referências, seguir as normas da ABNT, conforme o exemplo constante das referências. Após o título da seção haverá um espaço (linha em branco) simples.\\
    As considerações finais deverão apresentar os resultados do estudo e não deverá conter citações. \\



\section*{\large REFERÊNCIAS}


\begin{flushleft}

SOBRENOME, Nome. \textbf{Título da obra em negrito}: subtítulo sem negrito. Cidade: Editora, Ano. 

\

SOBRENOME, Nome.\textbf{Título da obra em negrito}. Cidade: Editora, Ano.

\

Poderão conter referências de internet desde que seguidas as normas da ABNT, 14ª edição. 



\end{flushleft}
}
\end{document}
